% DIN：给定候选对历史行为做局部激活（候选相关注意力），再加权求和，与其他字段拼接送入CTR头。
% 跨栏宽图（figure*），约150mm 宽。
\begin{tikzpicture}[x=1mm,y=1mm,
  font=\figurefont\footnotesize,
  dbox/.style={rectangle,draw=FigureInput!70,thick,fill=FigureInput!10,
    align=center,text=BookInk,rounded corners=1.4pt,minimum height=8mm,
    text width=22mm,inner sep=2pt},
  dhist/.style={rectangle,draw=FigureModel!70,fill=FigureModel!10,
    align=center,text=BookInk,rounded corners=1.4pt,minimum height=7.5mm,
    text width=20mm,inner sep=2pt},
  dact/.style={rectangle,draw=FigureLoss!75,fill=FigureLoss!14,
    align=center,text=BookInk,rounded corners=1.4pt,minimum height=8.5mm,
    text width=22mm,inner sep=2pt},
  dout/.style={rectangle,draw=FigureOutput!70,thick,fill=FigureOutput!12,
    align=center,text=BookInk,rounded corners=1.4pt,minimum height=10mm,
    text width=24mm,inner sep=2pt},
  dflow/.style={-{Stealth[length=1.8mm,width=1.2mm]},draw=FigureOutput!85,
    line width=.85pt},
  dfeed/.style={-{Stealth[length=1.8mm,width=1.2mm]},draw=FigureInput!85,
    line width=.9pt},
  dlab/.style={font=\figurefont\footnotesize,text=BookMuted,fill=white,inner sep=1pt}
]
  % 历史行为
  \node[dhist] (b1) at (14,18) {行为1 $i_1$\\背包};
  \node[dhist] (b2) at (14,9)  {行为2 $i_2$\\水壶};
  \node[dhist] (b3) at (14,0)  {行为3 $i_3$\\相机};
  \node[dhist] (b4) at (14,-9) {$\vdots$ 行为 $n$};
  \node[align=center,text=BookMuted] at (14,-15)
    {历史 $H_{u,t}$，嵌入 $\bm v_{i_k}$};

  % 当前候选与注入总线
  \node[dbox,fill=FigureInput!18] (cand) at (14,-25) {当前候选 $i$，$\bm v_i$};
  \coordinate (bus) at (37,-16);
  \draw[dfeed] (cand.north) -- (37,-21) -- (bus);
  \draw[dfeed] (bus) -- (37,22);
  \node[dlab,align=left,anchor=west] at (39,3) {候选 $\bm v_i$ 注入\\每个激活单元};

  % 行为 + 候选 -> 局部激活
  \node[dact] (a1) at (66,18) {局部激活\\$a_1(i)$};
  \node[dact] (a2) at (66,9)  {局部激活\\$a_2(i)$};
  \node[dact] (a3) at (66,0)  {局部激活\\$a_3(i)$};
  \node[dact] (a4) at (66,-9) {局部激活\\$a_n(i)$};
  \foreach \y in {18,9,0,-9}{
    \draw[dflow] (25,\y) -- (54,\y);
  }

  % 加权求和
  \node[dout,fill=FigureOutput!16,text width=30mm] (sum) at (100,5)
    {加权求和\\$\bm h(i)=\sum_k a_k(i)\bm v_{i_k}$};
  \foreach \y in {18,9,0,-9}{
    \draw[dflow] (78,\y) -- (88,\y) -- (88,5) -- (sum.west);
  }

  % 其他字段与预测头
  \node[dbox,fill=white,text width=26mm] (other) at (100,-13)
    {用户、情境等字段 $\bm x_t$};
  \draw[dflow] (other.north) -- (sum.south);
  \node[dout,fill=FigureOutput!20,text width=26mm] (mlp) at (134,5)
    {CTR 预测网络\\$\widehat p_i=\sigma(f_\theta)$};
  \draw[dflow] (sum.east) -- (mlp.west);
  \draw[dflow] (other.east) -- (127,-13) -- (127,-2) -- (mlp.south);
\end{tikzpicture}
